% ID: SEL-D0056-OA-E01
% ESTADO: validada. S057.
\textbf{a)} Sumamos $CD$ y distribuimos $(I_3+D)C$. Como $C$ y $D$ son simétricas, $DC=(CD)^t$.
\[
CD=\begin{pmatrix}1&0&1\\1&2&1\\1&0&1\end{pmatrix},
\qquad
X=\tfrac12(C+DC+CD)
=\boxed{\begin{pmatrix}1&1&1\\1&2&1\\1&1&1\end{pmatrix}}.
\]
\textbf{b)} Un $1$ fuera de la diagonal indica una arista. La simetría implica que las aristas son no dirigidas y las diagonales nulas excluyen lazos. $C$ da el camino $a-b-c$ y $D$ el triángulo de vértices $1,2,3$.
\begin{center}
\setlength{\unitlength}{0.9mm}
\begin{picture}(113,35)(-5,-5)
\qbezier(5,7)(25,7)(45,7)
\put(5,7){\makebox(0,0){\rule{1.4pt}{1.4pt}}}\put(25,7){\makebox(0,0){\rule{1.4pt}{1.4pt}}}\put(45,7){\makebox(0,0){\rule{1.4pt}{1.4pt}}}
\put(4,1){$a$}\put(24,1){$b$}\put(44,1){$c$}
\qbezier(66,7)(76,16)(86,25)
\qbezier(86,25)(96,16)(106,7)
\qbezier(106,7)(86,7)(66,7)
\put(66,7){\makebox(0,0){\rule{1.4pt}{1.4pt}}}\put(86,25){\makebox(0,0){\rule{1.4pt}{1.4pt}}}\put(106,7){\makebox(0,0){\rule{1.4pt}{1.4pt}}}
\put(65,1){$1$}\put(85,28){$2$}\put(105,1){$3$}
\end{picture}
\end{center}
